\section{El escenario final: entre universos, entre mundos, entre cuerpos}

Esta sección analiza el modelo del escenario final, más lejano. Xie Chen (谢晨) considera que el modelo final debe ser cross universe, cross world, cross embodiment, es decir, capaz de transferirse entre universos, entre mundos y entre cuerpos. La formulación suena grandiosa, pero detrás hay un objetivo sencillo: mejorar la generalización. Los seres humanos pueden pasar con rapidez de la realidad a un juego, de la cocina a la oficina, de la mano a una herramienta; si un robot solo puede trabajar en una simulación, en una habitación o con un cuerpo, todavía no es inteligencia general.

Los datos de videojuegos tienen aquí un significado especial. Los juegos no son necesariamente fieles a la física, pero ofrecen mundos ricos, reglas, objetivos e interacción en primera persona, por lo que resultan adecuados para mejorar, en la etapa de preentrenamiento, la capacidad del agent de transferirse entre mundos. DeepMind y otras instituciones usan grandes cantidades de datos de juegos precisamente porque el entrenamiento entre universos acostumbra al modelo a actuar con reglas y distribuciones de observación distintas. No sustituye el Real2Sim físico, pero complementa una world diversity más amplia.

\begin{figure}[H]
\centering
\includegraphics[width=0.92\textwidth]{cross-universe-model.png}
\caption{Entre universos, entre mundos, entre cuerpos: el modelo del escenario final debe mejorar la generalización entre entornos. Diagrama conceptual propio, elaborado a partir de la conversación de 01:21:25--01:23:28.}
\end{figure}

\begin{knowledgebox}{Lectura de la figura: cada tipo de transferencia resuelve un problema de generalización distinto}
La transferencia entre universos resuelve el domain shift entre simulación, juegos y realidad; la transferencia entre mundos resuelve las diferencias entre entornos como hogares, fábricas, hoteles y restaurantes; la transferencia entre cuerpos resuelve las diferencias entre los cuerpos de los robots y sus espacios de acción. Al leer la figura, hay que notar que las tres no son consignas, sino tres tipos de diferencias en la distribución de los datos. El modelo del escenario final debe ser capaz de extraer la esencia de la tarea a través de estas distribuciones.
\end{knowledgebox}

\begin{knowledgebox}{Términos clave: los tres tipos de transferencia}
\begin{tabularx}{\textwidth}{p{0.18\textwidth}p{0.36\textwidth}X}
\toprule
Capacidad & Significado & Requisitos de entrenamiento \\
\midrule
Entre universos & Transferirse entre juegos, simulación y el mundo real & Datos de mundos diversos, aprendizaje de reglas y domain adaptation. \\
Entre mundos & Transferirse entre distintos entornos de tarea & Los recursos de escena, la distribución de tareas y la retroalimentación real deben ser suficientemente amplios. \\
Entre cuerpos & Transferirse entre distintos cuerpos de robot & Hay que modelar el espacio de acción, los sensores, la geometría del cuerpo y las restricciones de control. \\
\bottomrule
\end{tabularx}
\end{knowledgebox}

\subsection{Todavía en la etapa GPT-1}

Esta subsección vuelve al diagnóstico de la etapa actual. Xie Chen considera que la inteligencia corporizada en su conjunto sigue en la etapa GPT-1; esto no significa que no tenga ninguna capacidad, sino que aún no se ha encontrado una receta estable de scaling law. La analogía con el FSD de Tesla resulta útil: al principio, añadir datos de forma continua no siempre producía mejoras sostenidas, hasta que se consolidó la vía de extremo a extremo; solo entonces la expansión de datos, cómputo y modelo empezó a traducirse de manera más estable en mejoras de capacidad. La inteligencia corporizada todavía no ha vivido ese momento claro.

Sin embargo, no es pesimista, por tres razones: la densidad de equipos fundadores y de científicos que hoy entran en la inteligencia corporizada es mucho mayor que en los inicios de la conducción autónoma; la atención del capital y de la industria es mayor; y los modelos grandes y la conducción autónoma ya aportaron experiencia en scaling, transformer, volante de datos, RL y entrenamiento de extremo a extremo. Dicho de otro modo, es una etapa temprana, pero no una etapa temprana que empiece desde cero.

\begin{figure}[H]
\centering
\includegraphics[width=0.96\textwidth]{embodied-gpt1-stage.png}
\caption{La inteligencia corporizada sigue en la etapa GPT-1: aún no se ha encontrado una scaling law estable y la recipe de datos sigue en exploración. Diagrama conceptual propio, elaborado a partir de la conversación de 01:23:28--01:28:21.}
\end{figure}

\begin{knowledgebox}{Lectura de la figura: la etapa GPT-1 significa que la recipe no ha convergido}
Lo más importante de la figura es «unknown recipe». La industria sabe que necesita datos, cómputo, modelos, simulación y RL, pero todavía no sabe qué combinación mejorará de forma estable la capacidad de los robots reales. Más que indicar que la capacidad es débil, la etapa GPT-1 indica que aún no se ha formado una vía de escalamiento tan clara como la de GPT-3/ChatGPT.
\end{knowledgebox}

\begin{importantbox}{Scaling law moment}
El scaling law moment de la inteligencia corporizada no es una ronda de financiamiento, una demo ni el lanzamiento de un modelo aislado, sino el momento en que el equipo descubre que, con una arquitectura y una canalización de datos dadas, seguir añadiendo datos de alta calidad, cómputo y entornos mejora de forma predecible la capacidad en tareas reales.
\end{importantbox}

\subsection{El progress es más sano que el financiamiento}

Esta subsección conserva el juicio de emprendedor de la segunda mitad de la entrevista. A Xie Chen le preocupa que la industria se fije demasiado en los montos de financiamiento y en las valuaciones, y descuide el progress real. Desde la perspectiva de RL, si el reward model de un equipo se convierte en «dar conferencias, hacer demos y atraer inversionistas», su conducta queda guiada por una recompensa equivocada; un reward model más sano consiste en atender a los clientes, crear valor, cobrar y mejorar el producto de forma continua con la retroalimentación de los clientes.

Este pasaje es, en realidad, coherente con la línea técnica principal: ya se trate de datos, de simulación o de una empresa, todo necesita retroalimentación real. Para el modelo, es la retroalimentación del despliegue real; para la empresa, la retroalimentación de los pagos y las recompras de los clientes; para la industria, el progress real y no el narrative progress. En la entrevista, el paso de Guanglun (光轮) de la conducción autónoma a la inteligencia corporizada también implicó pasar de «yo sé de simulación» a «volver a aprender la inteligencia corporizada desde el nivel de licenciatura», lo que muestra que la experiencia de problemas de baja dimensión no puede aplicarse directamente a problemas de alta dimensión.

\begin{warningbox}{Una señal de burbuja en la industria}
Si un equipo optimiza sobre todo la narrativa para el financiamiento, la impresión que causan las demos y la valuación, en lugar del valor para el cliente, la mejora del modelo y el despliegue real, su reward model está desviado. En una industria con tanta carga de ingeniería como la inteligencia corporizada, una recompensa equivocada es más peligrosa que un fracaso técnico de corto plazo.
\end{warningbox}

\subsection{Resumen de la sección}

El escenario final de la inteligencia corporizada es la generalización entre universos, entre mundos y entre cuerpos; en la etapa actual todavía se busca la scaling law. Quizá la industria encuentre pronto la recipe, pero a condición de alinear el reward con el progress real: mejoras reales del modelo, valor real para el cliente y retroalimentación real del despliegue.
